\chapter{Machine settings by driver}
\label{app:drivers}

A \emph{driver} is the part of MeerK40t that speaks one family of laser
controller. Six families ship with the program: Lihuiyu (the K40 nano boards),
GRBL, Ruida, Moshi, Newly and Balor (the JCZ galvo controllers). Each driver
defines its own settings, and a device shows only the settings of its own
driver --- two machines with similar-looking Configuration windows can hold
quite different contents. This appendix is the lookup table for those settings.
Every row gives the label that appears in the device's \btn{Config} window, the
internal \emph{key}, the value the program ships with, and what the setting
changes.

Open the window from the \panel{Devices} pane with \btn{Config} for the active
device, or through \menu{Device-Settings>Device Manager}. The key column is the
name of the underlying variable, which is what the console \cmd{set} command
reads and writes (Section~\ref{ch:console}, Appendix~\ref{app:console}).

\begin{notebox}[Bed size and origin belong to the machine]
Width, Height and \ui{Force Declared Home} describe the hardware, not the
artwork. If you change them after placing a design, the job is re-mapped onto
the new bed and the same drawing may end up somewhere else on the material. Set
the bed up before you position anything, and check placement again afterwards
(Chapter~\ref{ch:where}).
\end{notebox}

\section*{Common device settings}

Most drivers offer the same handful of settings, usually under the same keys.
Those are listed once below. The driver sections that follow give the values
that differ, and the settings that only one family has. A driver is expected to
show only the rows relevant to it.

{\small
\begin{longtable}{@{}>{\raggedright\arraybackslash}p{3.2cm}@{\hspace{6pt}}>{\footnotesize\raggedright\arraybackslash}p{4.6cm}@{\hspace{6pt}}>{\footnotesize\raggedright\arraybackslash}p{2.6cm}@{\hspace{6pt}}>{\raggedright\arraybackslash}p{4.6cm}@{}}
\toprule
\textbf{Setting} & \textbf{Key} & \textbf{Default} & \textbf{What it changes} \\
\midrule
\endfirsthead
\toprule
\textbf{Setting} & \textbf{Key} & \textbf{Default} & \textbf{What it changes} \\
\midrule
\endhead
\midrule
\multicolumn{4}{r@{}}{\footnotesize\itshape continued overleaf} \\
\endfoot
\bottomrule
\endlastfoot
Label & \setting{label} & shipped name & Device name in the \panel{Devices} list. \\
Width, Height & \setting{bedwidth}, \setting{bedheight} & see each driver & Bed size. Sets the scene, the bed limits and jog confinement. \\
Laserspot & \setting{laserspot} & 0.3\unit{mm} & Beam spot size used for display and planning. All six drivers. \\
X-Axis, Y-Axis & \setting{scale_x}, \setting{scale_y} & 1.000 & Scale from board units to millimetres. Lihuiyu calls these \setting{user_scale_x} and \setting{user_scale_y}. \\
Flip X, Flip Y & \setting{flip_x}, \setting{flip_y} & off & Mirror an axis. On for \setting{flip_y} on GRBL and Balor. \\
Swap X and Y & \setting{swap_xy} & off & Exchanges the axes. Applied before the flips. On by default for Balor. \\
Force Declared Home & \setting{home_corner} & \setting{auto} & Which corner counts as home. All drivers except Balor. \\
Curve Interpolation & \setting{interp}; Ruida: \setting{interpolate} & 5 mils; Ruida 500\textmu m & Distance at which curves are flattened into straight segments. \\
Device Position & \setting{signal_updates} & on & Draws the live position indicator. \\
Coolant & \setting{device_coolant} & nothing & Coolant or air-assist method claimed for the device. All except Newly. \\
Max vector speed, Max raster speed & \setting{max_vector_speed}, \setting{max_raster_speed} & 140, 750 & Highest reliable cutting and raster speeds. Lihuiyu and GRBL. \\
Travel speed & \setting{rapid_speed} (Balor: \setting{default_rapid_speed}) & see drivers & Speed used for moves that do not cut. \\
Operation Defaults, speed & \setting{default_speed_op_cut}, \setting{default_speed_op_engrave} and the matching keys for raster and image & Cut 5, Engrave 25, Raster 250 & Speed given to a newly created operation. Balor uses 150, 250 and 500. \\
Operation Defaults, power & \setting{default_power_op_cut} and the same key pattern for the other three & 1000 & Power given to a newly created operation, in whatever unit the row below selects. \\
Show power as \% & \setting{use_percent_}\allowbreak\setting{for_power_display} & on for GRBL, Newly, Balor & On: full power reads as 100\%. Off: full power reads as 1000. \\
Show speed in mm/min & \setting{use_mm_min_for_}\allowbreak\setting{speed_display} & off & On: speeds are shown in mm/min. Off: mm/s. The Lihuiyu, Ruida and Moshi sources keep the same preference under the parallel key \setting{use_minute_}\allowbreak\setting{for_speed_display}. \\
Limit the controller buffer size, Controller Buffer & \setting{limit_buffer}, \setting{max_buffer} & on, 200 lines & GRBL: stop sending once the controller holds the stated number of lines. \\
Write Buffer & \setting{buffer_max}, \setting{buffer_limit} & 900, on & Lihuiyu: cap on the USB write queue. The Interface panel exposes it as \ui{Limit Write Buffer} with its own maximum. \\
X-Margin, Y-Margin & \setting{user_margin_x}, \setting{user_margin_y} & 0 & Present in the source but marked as not working, so no driver shows it in the window. \\
\end{longtable}
}

Two things that look like driver settings but are not. Homing on start is a job
step, not a device setting: add \ui{Append Physical Home} or
\ui{Append Return to Origin} in the Operations tree (Chapter~\ref{ch:order}).
And the coolant list only expands after a coolant method has been registered, so
the combo may hold nothing but \ui{Nothing} until you use one.

\section*{Lihuiyu (K40)}

The Lihuiyu driver runs the M2 and M3 nano boards found in K40-style CO\textsubscript{2}
cutters and their clones. The board is a USB device reached through the CH341
chip: choose \ui{USB} in the Interface panel for the normal case, \ui{Networked}
to send the same packets to another MeerK40t instance over TCP, or \ui{Mock} for
a dry run with no machine attached. These boards have no endstops, so \ui{Home}
moves to the declared origin rather than seeking a switch. The Configuration
window carries the tabs Configuration, Interface, Setup, Effects, Operation
Defaults, Warning, Default Actions and Display Options.

{\small
\begin{longtable}{@{}>{\raggedright\arraybackslash}p{3.2cm}@{\hspace{6pt}}>{\footnotesize\raggedright\arraybackslash}p{4.6cm}@{\hspace{6pt}}>{\footnotesize\raggedright\arraybackslash}p{2.6cm}@{\hspace{6pt}}>{\raggedright\arraybackslash}p{4.6cm}@{}}
\toprule
\textbf{Setting} & \textbf{Key} & \textbf{Default} & \textbf{What it changes} \\
\midrule
\endfirsthead
\toprule
\textbf{Setting} & \textbf{Key} & \textbf{Default} & \textbf{What it changes} \\
\midrule
\endhead
\midrule
\multicolumn{4}{r@{}}{\footnotesize\itshape continued overleaf} \\
\endfoot
\bottomrule
\endlastfoot
Device Name & \setting{label} & \setting{M2-Nano}, \setting{M3-Nano} & Internal name used for the device; \setting{lihuiyu-device} when a device is added without a model entry. \\
Width, Height & \setting{bedwidth}, \setting{bedheight} & 310\unit{mm}, 210\unit{mm} & Bed size of the usual K40. \\
X Scale Factor, Y Scale Factor & \setting{user_scale_x}, \setting{user_scale_y} & 1.000 & Board units to millimetres. \\
Board & \setting{board} & \setting{M2} & Which nano board is fitted: \setting{M2}, \setting{M3}, \setting{B2}, \setting{M}, \setting{M1}, \setting{A}, \setting{B} or \setting{B1}. This changes the speed codes sent to the board. \\
Hardware-PWM & \setting{supports_pwm} & off & Power control by hardware PWM. Only M3 boards with firmware 2024.01.18g or newer; earlier M3 revisions behave as M2+. Shown only for the M3 board. \\
Maximum Laser strength & \setting{power_scale} & 30\% & Ceiling that every operation's power is a fraction of; adjustable from 1 to 50, or to 100 with developer mode. Setting it too high can damage the laser tube. \\
Strict & \setting{strict} & off & Enter and leave speed mode travelling in the same direction (an M2-V4 fix). \\
Twitch Vectors & \setting{twitches} & off & Re-enables the old twitch behaviour when moving the vectors. \\
Use legacy raster method & \setting{legacy_raster} & on & Uses the legacy raster engine: better at high speed, but with artifacts. \\
Rapid Moves Between Objects & \setting{opt_rapid_between} & on & Travel at rapid speed between separate objects. \\
Minimum Jog Distance & \setting{opt_jog_minimum} & 256 & Step below which a move is treated as a jog. \\
Jog Method & \setting{opt_jog_mode} & 0 & \ui{Default}, \ui{Reset} or \ui{Finish}. \\
Override Rapid Movements & \setting{rapid_override} & off & Replaces the travel speed with the two values below. \\
X Travel Speed, Y Travel Speed & \setting{rapid_override_speed_x}, \setting{rapid_override_speed_y} & 50.0 & Per-axis travel speed used while the override is on. \\
Max vector speed, Max raster speed & \setting{max_vector_speed}, \setting{max_raster_speed} & 140, 750 & Highest reliable speeds. \\
Automatically lock rail & \setting{autolock} & on & Locks the rail when a job finishes. \\
Phase Type, Phase Value & \setting{plot_phase_type}, \setting{plot_phase_value} & 0, 0 & Raster phase grouping: \ui{Sequential}, \ui{Random}, \ui{Progressive} or \ui{Static}, with a value out of 1000. \\
Pulse Grouping & \setting{plot_shift} & off & Groups pulses in raster output. \\
Fix rated to actual speed & \setting{fix_speeds} & off & These boards run at about 92\% of the rated speed; enable this to correct the displayed speed. \\
Networked address, port & \setting{address}, \setting{port} & localhost, 1025 & Host and port used when the interface is set to \ui{Networked}. Windows builds default to port 1022. \\
Write Buffer & \setting{buffer_max}, \setting{buffer_limit} & 900, on & Cap on the USB write queue, with the Interface panel's \ui{Limit Write Buffer} as the visible control. \\
\end{longtable}
}

Outside the configuration window, the acceleration chart holds the vector and
raster acceleration breakpoints: the groups \ui{Vector} and \ui{Vertical Raster}
default to 25.4, 60, 127 and infinity, \ui{Horizontal Raster} to 25.4, 127, 320
and infinity, each group with its own enable checkbox.

\section*{GRBL}

The GRBL driver covers a large group of diode and CO\textsubscript{2} machines:
Ortur and Longer diode engravers, generic GRBL boards, FluidNC controllers and
K40 cutters fitted with a GRBL controller such as the MKS DLC32. It reaches the
board over a serial port (a USB-serial adapter appears as a port), over TCP or
WebSocket for boards with a network socket, or over a mock. The shipped device
entries already fill in the bed size and homing details for common machines, so
start from the entry that matches your machine. The Configuration window has the
tabs Device, Interface, Protocol, Advanced, Effects, Operation Defaults, ESP3D
Upload, Warning, Default Actions and Display Options.

{\small
\begin{longtable}{@{}>{\raggedright\arraybackslash}p{3.2cm}@{\hspace{6pt}}>{\footnotesize\raggedright\arraybackslash}p{4.6cm}@{\hspace{6pt}}>{\footnotesize\raggedright\arraybackslash}p{2.6cm}@{\hspace{6pt}}>{\raggedright\arraybackslash}p{4.6cm}@{}}
\toprule
\textbf{Setting} & \textbf{Key} & \textbf{Default} & \textbf{What it changes} \\
\midrule
\endfirsthead
\toprule
\textbf{Setting} & \textbf{Key} & \textbf{Default} & \textbf{What it changes} \\
\midrule
\endhead
\midrule
\multicolumn{4}{r@{}}{\footnotesize\itshape continued overleaf} \\
\endfoot
\bottomrule
\endlastfoot
Label & \setting{label} & \setting{grbl} & Device name. The shipped entries use names such as \setting{GRBL-DLC32-400} or \setting{GRBL-K40-CO2}. \\
Width, Height & \setting{bedwidth}, \setting{bedheight} & 235\texttimes235\unit{mm} & Bed size. The entries carry their own: 405\texttimes285\unit{mm} for the MKS DLC32, 400\texttimes430\unit{mm} for the Ortur, 450\texttimes450\unit{mm} for the Longer Ray5. \\
X-Axis, Y-Axis & \setting{scale_x}, \setting{scale_y} & 1.000 & Board units to millimetres. \\
Flip X & \setting{flip_x} & off & Mirrors X. +X is standard for GRBL; leave it off unless your machine differs. \\
Flip Y & \setting{flip_y} & on & On this machine Y goes negative into the bed. Leave it on with a top-left home corner, or Y jogs can trip the soft limit (ALARM:2). \\
Swap XY & \setting{swap_xy} & off & Exchanges the axes before the flips. \\
Force Declared Home & \setting{home_corner} & \setting{auto} & Set the corner here if the machine homes somewhere other than the default. \\
Supports Z-axis & \setting{supports_z_axis} & off & Enables Z commands for a machine with a Z axis. \\
Z-Homing & \setting{z_home_command} & \cmd{$HZ} & Which command starts the Z homing sequence: \cmd{$HZ} or \cmd{G28 Z}. \\
Interface Type & \setting{interface} & \setting{serial} & \ui{Serial}, \ui{TCP-Network}, \ui{WebSocket-Network} or \ui{Mock}. \\
(serial port) & \setting{serial_port} & \cmd{UNCONFIGURED} & Combo of the ports detected on the computer. \\
Baud Rate & \setting{baud_rate} & 115200 & Serial speed. \\
Address, Port & \setting{address}, \setting{port} & localhost, 23 & Host and port for the TCP interface. \\
Address, Port (WebSocket) & \setting{address}, \setting{port} & localhost, 81 & Host and port for the WebSocket interface. \\
Use M3 & \setting{use_m3} & off & Starts the laser with M3 instead of M4. \\
Use G1 before Rapid Moves & \setting{use_g1_for_power} & off & Sends \cmd{G1 S0} instead of \cmd{G0 X Y S0}, for firmwares that mishandle rapid moves carrying power. \\
Reset on 'Clear Alarm' & \setting{extended_alarm_clear} & off & Clearing an alarm also soft-resets the controller. \\
Device has endstops & \setting{has_endstops} & off & Enables physical homing with \cmd{$H}. \\
Sequential homing (\cmd{$HY} then \cmd{$HX}) & \setting{sequential_homing} & off & For the MKS DLC32 and machines whose limit switches sit at the top left. Homing both axes at once can stop about 50\unit{mm} short with an ALARM:1. \\
Simulate reddot & \setting{use_red_dot} & off & Burns a low-power visible dot for framing. GRBL has no separate red-dot output pin; this uses the laser itself. \\
Reddot Laser strength & \setting{red_dot_level} & 30 & Strength of that dot, out of 1000. \\
Active during Outline & \setting{use_red_dot_for_outline} & on & Uses the red dot automatically when outlining the burn area. \\
Max vector speed & \setting{max_vector_speed} & 140 & Highest reliable speed for cutting and engraving moves. \\
Max raster speed & \setting{max_raster_speed} & 750 & Highest reliable speed for raster moves. \\
Travel speed & \setting{rapid_speed} & 600 & Speed of moves that do not cut. \\
Limit the controller buffer size & \setting{limit_buffer} & on & Enables the buffer limit below. \\
Controller Buffer & \setting{max_buffer} & 200 & Lines the controller may hold before sending pauses. \\
Require Validator & \setting{require_validator} & off & Waits for the expected welcome text before validating a job. \\
Welcome Validator & \setting{welcome} & \cmd{Grbl} & The welcome text to wait for. \\
Reset on connect & \setting{reset_on_connect} & off & Soft-resets the controller when connecting. \\
Check sequence on connect. & \setting{boot_connect_sequence} & on & Runs the information query when connecting. \\
Sending Protocol & \setting{buffer_mode} & \setting{buffered} & \ui{buffered} or \ui{sync}. \\
Planning Buffer Size & \setting{planning_buffer_size} & 128 & Size of the planning queue. \\
Line Ending & \setting{line_end} & \setting{LF} & \ui{CR}, \ui{LF} or \ui{CRLF}. LF is required for the SD-card files of the MKS DLC32. \\
Post Connection Delay & \setting{connect_delay} & 0 & Milliseconds to wait after connecting before sending. \\
Startup commands & \setting{startup_commands} & empty & Multiline G-code sent on connect. \\
Enable ESP3D Upload & \setting{esp3d_enabled} & off & Shows the ESP3D buttons on the ribbon and unlocks the page below. \\
ESP3D Host, Port, SD Path & \setting{esp3d_host}, \setting{esp3d_port}, \setting{esp3d_path} & 192.168.1.100, 80, \file{/} & Where the ESP3D web interface lives and where its SD card is rooted. \\
ESP3D Username, Password & \setting{esp3d_username}, \setting{esp3d_password} & empty, empty & Optional login for that interface. \\
Cleanup Local Files & \setting{esp3d_cleanup} & on & Deletes the local G-code once it has been uploaded. \\
Prepare Plan exports for SD card & \setting{sd_export_prepare} & on & Patches a Plan export for a MKS DLC32 SD card: LF line endings, M3 laser mode, no \cmd{G28}. Home the machine yourself before running such a file. \\
\end{longtable}
}

\section*{Ruida}

The Ruida driver targets CO\textsubscript{2} machines with a Ruida controller,
the K50/K60 class. It connects over USB serial --- the Serial Interface field
takes a Linux device path or a Windows \cmd{COM} port --- or over the network,
where the driver talks UDP to the controller. It also ships the
\cmd{ruidacontrol} emulator, which lets MeerK40t answer as a Ruida controller
for LightBurn on UDP 50200 (jog on 50207), with optional man-in-the-middle
(40200/40207) and LB2RD bridge (TCP 5005) modes. The Configuration window has
the tabs Ruida, Interface, Configuration, Effects, Operation Defaults, Warning,
Default Actions and Display Options.

\begin{warnbox}[Under revision]
The Ruida module is marked incomplete in the project. It has been tested on
Linux only; the network emulation is described as untested and direct hardware
control as not implemented. The project's own notes also require raster layers
to use overscan 0 and, on the controller that was tested, Flip X to be enabled.
Treat this driver as experimental.
\end{warnbox}

{\small
\begin{longtable}{@{}>{\raggedright\arraybackslash}p{3.2cm}@{\hspace{6pt}}>{\footnotesize\raggedright\arraybackslash}p{4.6cm}@{\hspace{6pt}}>{\footnotesize\raggedright\arraybackslash}p{2.6cm}@{\hspace{6pt}}>{\raggedright\arraybackslash}p{4.6cm}@{}}
\toprule
\textbf{Setting} & \textbf{Key} & \textbf{Default} & \textbf{What it changes} \\
\midrule
\endfirsthead
\toprule
\textbf{Setting} & \textbf{Key} & \textbf{Default} & \textbf{What it changes} \\
\midrule
\endhead
\midrule
\multicolumn{4}{r@{}}{\footnotesize\itshape continued overleaf} \\
\endfoot
\bottomrule
\endlastfoot
Label & \setting{label} & \setting{ruida} & Device name. \\
Width, Height & \setting{bedwidth}, \setting{bedheight} & 24\unit{in}, 16\unit{in} & Bed size. \\
X-Axis, Y-Axis & \setting{scale_x}, \setting{scale_y} & 1.000 & Board units to millimetres. \\
Flip X & \setting{flip_x} & off & Mirrors X. The project's notes require this to be on for the tested machine. \\
Flip Y & \setting{flip_y} & off & Mirrors Y. \\
Swap X and Y & \setting{swap_xy} & off & Exchanges the axes. \\
Force Declared Home & \setting{home_corner} & \setting{auto} & Which corner counts as home. \\
Curve Interpolation & \setting{interpolate} & 500\textmu m & Distance at which curves are flattened. \\
Laser Power & \setting{default_power} & 20.0 & Global default power, overridden by each operation. \\
Cut Speed & \setting{default_speed} & 40.0 & Global default speed, overridden by each operation. \\
Swizzle Magic Number & \setting{magic} & \cmd{0x88} & Value used to encode the job file sent to the laser. \\
Baud Rate & \setting{baud_rate} & 115200 & Serial speed. \\
Serial Interface & \setting{serial} & empty & Serial port the controller is on, for example a Linux device path or \cmd{COM3}. \\
Interface & \setting{interface} & \setting{usb} & \ui{USB} or \ui{Networked} (UDP). The address field defaults to localhost and is only used when networked. \\
Coolant & \setting{device_coolant} & empty & Coolant or air-assist method. \\
\end{longtable}
}

Ruida jobs are exported as Ruida job files with the \file{.rd} extension. The
Plan export uses \cmd{save_job <file>} and applies the swizzle magic number
unless another value is passed for it. The ribbon adds \btn{Focus Z}, which
sends the controller's FocusZ command.

\section*{Moshi}

The Moshi driver covers CO\textsubscript{2} cutters with a Moshiboard, a clone
family that presents itself to the computer much like a Lihuiyu board. It
connects over USB through the same CH341 chip and offers no network interface.
There are no endstops, so \ui{Home} is simply a move to position 0,0; where the
job sits on the material is decided entirely by the bed size and the home
corner. The Configuration window has the tabs Configuration, Effects, Operation
Defaults, Warning, Default Actions and Display Options.

{\small
\begin{longtable}{@{}>{\raggedright\arraybackslash}p{3.2cm}@{\hspace{6pt}}>{\footnotesize\raggedright\arraybackslash}p{4.6cm}@{\hspace{6pt}}>{\footnotesize\raggedright\arraybackslash}p{2.6cm}@{\hspace{6pt}}>{\raggedright\arraybackslash}p{4.6cm}@{}}
\toprule
\textbf{Setting} & \textbf{Key} & \textbf{Default} & \textbf{What it changes} \\
\midrule
\endfirsthead
\toprule
\textbf{Setting} & \textbf{Key} & \textbf{Default} & \textbf{What it changes} \\
\midrule
\endhead
\midrule
\multicolumn{4}{r@{}}{\footnotesize\itshape continued overleaf} \\
\endfoot
\bottomrule
\endlastfoot
Label & \setting{label} & \setting{Moshiboard} & Device name. \\
Width, Height & \setting{bedwidth}, \setting{bedheight} & 330\unit{mm}, 210\unit{mm} & Bed size. \\
X-Axis, Y-Axis & \setting{scale_x}, \setting{scale_y} & 1.000 & Board units to millimetres. \\
Flip X, Flip Y & \setting{flip_x}, \setting{flip_y} & off, off & Mirror an axis. \\
Swap XY & \setting{swap_xy} & off & Exchanges the axes before the flips. \\
Force Declared Home & \setting{home_corner} & \setting{auto} & Which corner counts as home. \\
Curve Interpolation & \setting{interp} & 5 & Distance at which curves are flattened, in mils. \\
Use legacy raster method & \setting{legacy_raster} & on & Legacy raster engine: better at high speed, but with artifacts. \\
Run mock-usb backend & \setting{mock} & off & Connects to a fake laser instead of the real board, for debugging. \\
Device Position & \setting{signal_updates} & on & Draws the live position indicator. \\
Coolant & \setting{device_coolant} & empty & Coolant or air-assist method. \\
Travel speed & \setting{rapid_speed} & 40 & Speed of moves that do not cut. Present in the source but not shown in the window. \\
\end{longtable}
}

\section*{Newly}

The Newly driver covers CO\textsubscript{2} machines whose controller presents
itself as a Newly or U-SET board: the RayLaser class and about twenty clones,
including the Beijing SZTaiming, Gama, HPC, Rabit, Lion and Jinan families. It
connects over USB and has no network option. Its distinguishing idea is the file
slot: a job is sent to one of ten files on the machine, File 0 runs immediately,
and the others can be sent and started later from the ribbon (\btn{File 0} to
\btn{File 9}, \btn{Send Only} or \btn{Send \& Start}, \btn{Draw Frame},
\btn{Move Frame}, \btn{Replay}). The Configuration window has the tabs Newly,
Device, Global, Raster Chart, Effects, Operation Defaults, Warning, Default
Actions and Display Options.

{\small
\begin{longtable}{@{}>{\raggedright\arraybackslash}p{3.2cm}@{\hspace{6pt}}>{\footnotesize\raggedright\arraybackslash}p{4.6cm}@{\hspace{6pt}}>{\footnotesize\raggedright\arraybackslash}p{2.6cm}@{\hspace{6pt}}>{\raggedright\arraybackslash}p{4.6cm}@{}}
\toprule
\textbf{Setting} & \textbf{Key} & \textbf{Default} & \textbf{What it changes} \\
\midrule
\endfirsthead
\toprule
\textbf{Setting} & \textbf{Key} & \textbf{Default} & \textbf{What it changes} \\
\midrule
\endhead
\midrule
\multicolumn{4}{r@{}}{\footnotesize\itshape continued overleaf} \\
\endfoot
\bottomrule
\endlastfoot
Label & \setting{label} & \setting{newly-device} & Device name. \\
Width, Height & \setting{bedwidth}, \setting{bedheight} & 310\unit{mm}, 210\unit{mm} & Bed size. \\
Output File & \setting{file_index} & 0 & Which machine file the job is sent to, 0 to 9. File 0 executes immediately; the rest must be sent and then started. \\
Automatically Start & \setting{autoplay} & on & Starts the job as soon as the output file has been sent. Uncheck to send without running. \\
Machine index to select & \setting{machine_index} & 0 & Which attached board to use; leave at 0 with a single machine. \\
Max Power & \setting{max_power} & 20.0\% & Ceiling that operation power is scaled against. \\
Max Pulse Power & \setting{max_pulse_power} & 65.0\% & Power level at which pulses are fired. \\
PWM Power & \setting{pwm_enabled} & off & Enables pulse-width modulation for power control. \\
PWM Frequency & \setting{pwm_frequency} & 2 kHz & 1, 2, 3, 4, 5, 10, 20, 50, 75, 100, 200 or 255 kilohertz. Shown only when PWM Power is on. \\
Horizontal DPI, Vertical DPI & \setting{h_dpi}, \setting{v_dpi} & 1000, 1000 & Dots per inch along each axis. \\
Horizontal Backlash, Vertical Backlash & \setting{h_backlash}, \setting{v_backlash} & 0\unit{mm}, 0\unit{mm} & Backlash compensation per axis. \\
Cut Speed, Cut Power & \setting{default_cut_speed}, \setting{default_cut_power} & 15.0, 1000.0 & Global defaults for cutting; operations override them. \\
Raster Speed & \setting{default_raster_speed} & 200.0 & Global default raster speed. \\
On Delay, Off Delay & \setting{default_on_delay}, \setting{default_off_delay} & 0\unit{ms}, 0\unit{ms} & Delay before the laser turns on and after it turns off. \\
Cut DC, Move DC & \setting{cut_dc}, \setting{move_dc} & 100, 100 & Current setting for cutting moves and for regular moves. \\
Maximum Raster Jog & \setting{max_raster_jog} & 15 & Longest distance allowed in a single raster step or jog. \\
Speed Chart & \setting{speedchart} & five rows & Raster chart of speed against acceleration length, backlash and corner speed: 100/8/0/20, 200/10/0/20, 300/14/0/20, 400/16/0/20 and 500/18/0/20. Rows can be added, changed and removed. \\
Run mock-usb backend & \setting{mock} & off & Connects to a fake laser instead of the real board. \\
Curve Interpolation & \setting{interp} & 5 & Distance at which curves are flattened, in mils. \\
Force Declared Home & \setting{home_corner} & \setting{auto} & Which corner counts as home. \\
Flip X, Flip Y & \setting{flip_x}, \setting{flip_y} & off, off & Mirror an axis. \\
Swap X and Y & \setting{swap_xy} & off & Exchanges the axes. \\
Device Position & \setting{signal_updates} & on & Draws the live position indicator. \\
\end{longtable}
}

\section*{Balor (galvo)}

The Balor driver runs galvo marking heads built on a JCZ controller: fibre,
CO\textsubscript{2} and UV sources, including MOPA fibre lasers. It connects
over USB only --- there is no host or port setting, and the device location
reports either \setting{usb} or \setting{mock} --- with a machine index and an
optional serial-number check to pick the right head when several are attached.
The bed is the lens field, so the Width setting is labelled \ui{Lens Size} and
defaults to 110\unit{mm}; a correction file can replace it with the measured
field. The Configuration window has the tabs Balor, Redlight, Global, Timings,
Extras, Effects, Operation Defaults, Warning, Default Actions and Display
Options.

{\small
\begin{longtable}{@{}>{\raggedright\arraybackslash}p{3.2cm}@{\hspace{6pt}}>{\footnotesize\raggedright\arraybackslash}p{4.6cm}@{\hspace{6pt}}>{\footnotesize\raggedright\arraybackslash}p{2.6cm}@{\hspace{6pt}}>{\raggedright\arraybackslash}p{4.6cm}@{}}
\toprule
\textbf{Setting} & \textbf{Key} & \textbf{Default} & \textbf{What it changes} \\
\midrule
\endfirsthead
\toprule
\textbf{Setting} & \textbf{Key} & \textbf{Default} & \textbf{What it changes} \\
\midrule
\endhead
\midrule
\multicolumn{4}{r@{}}{\footnotesize\itshape continued overleaf} \\
\endfoot
\bottomrule
\endlastfoot
Label & \setting{label} & \setting{balor-device} & Device name. \\
Laser Source & \setting{source} & \setting{fiber} & \ui{fiber}, \ui{co2} or \ui{uv}. \\
Enable (correction file) & \setting{corfile_enabled} & off & Uses a lens correction file. \\
File & \setting{corfile} & none & The \file{.cor} correction file. Loading one updates Lens Size from its scale. \\
Width (Lens Size) & \setting{lens_size} & 110\unit{mm} & Size of the galvo field; the bed is this square. \\
Invert Power & \setting{power_invert} & off & Inverts the power scale. Shown for the UV source. \\
Flip X, Flip Y & \setting{flip_x}, \setting{flip_y} & off, on & Mirror an axis. \\
Swap XY & \setting{swap_xy} & on & Exchanges the axes. \\
Rotate View & \setting{rotate} & 0 & Rotates the working view: 0, 90, 180 or 270 degrees. \\
Curve Interpolation & \setting{interp} & 5 & Distance at which curves are flattened. \\
Machine index to select & \setting{machine_index} & 0 & Which attached controller to use. \\
Check serial no & \setting{serial_enable} & off & Requires a specific head, matched by serial number. \\
serial & \setting{serial} & empty & The serial number to match. Shown when the check is enabled. \\
Footpedal & \setting{footpedal_pin} & 15 & GPIO pin the foot pedal is wired to. \\
Redlight laser & \setting{light_pin} & 8 & GPIO pin the red pilot laser is wired to. \\
Pedal action & \setting{pedal_mode} & \setting{ignore} & What a pedal press does: \setting{ignore}, pause and resume the job, pause while pressed, stop the job, or arm and start the job. \\
Active on low signal & \setting{pedal_active_low} & on & Treats the pedal as pressed when the signal is low. \\
Pedal poll interval & \setting{pedal_poll_interval} & 0.25 & Seconds between checks of the pedal. \\
Redlight travel speed & \setting{redlight_speed} & 3000 & Speed of the galvo while using the red pilot laser. \\
Dark, Light & \setting{redlight_delay_dark}, \setting{redlight_delay_light} & 1, 1 & Microseconds of delay for dark and for light redlight moves. \\
X-Offset, Y-Offset, Angle Offset & \setting{redlight_offset_x}, \setting{redlight_offset_y} & 0\unit{mm}, 0\unit{mm} & Shifts the redlight positions in x and in y. The angle offset, \setting{redlight_angle}, curves them around the centre and also defaults to 0\textdegree. \\
Prefer redlight on & \setting{redlight_preferred} & off & Turns a toggleable redlight back on when a job completes. \\
Restart light jobs & \setting{restart_light_jobs} & off & Restarts light jobs after they finish. \\
Power & \setting{default_power} & 500.0 & Global default power; operations override it. \\
Speed & \setting{default_speed} & 100.0 & Global default marking speed. \\
Q Switch Frequency & \setting{default_frequency} & 30.0 & Q-switch frequency in kilohertz. \\
First Pulse Killer & \setting{default_fpk} & 10.0\% & Percentage of first-pulse suppression, for the CO\textsubscript{2} and UV sources. \\
Travel Speed & \setting{default_rapid_speed} & 2000.0 & Speed used when not marking. \\
Enable (Pulse Width) & \setting{pulse_width_enabled} & off & MOPA heads: use pulse width as well as power. \\
Set Pulse Width (ns) & \setting{default_pulse_width} & 4 & Pulse width in nanoseconds, chosen from a fixed list: 1, 2, 4, 6, 9, 13, 20, 30, 45, 55, 60, 80, 100, 150, 200 or 250. Shown when pulse width is enabled. \\
Laser On, Laser Off & \setting{delay_laser_on}, \setting{delay_laser_off} & 100, 100 & Microseconds of delay at the start and at the end of each marking command. \\
Polygon Delay & \setting{delay_polygon} & 100 & Microseconds of delay between points along a path. \\
End Delay & \setting{delay_end} & 300 & Microseconds of delay at the end of the job. \\
Long jump delay, Short jump delay & \setting{delay_jump_long}, \setting{delay_jump_short} & 200, 8 & Microseconds of delay for long and for short jumps. \\
Run mock-usb backend & \setting{mock} & off & Connects to a fake controller instead of the real head. \\
Device Position & \setting{signal_updates} & on & Draws the live position indicator. \\
\end{longtable}
}

\section*{Where to check the current value}

The Configuration window always shows the live values, under the same labels as
the tables above: open the \panel{Devices} pane, select the device and press
\btn{Config}. Three console commands help when a window is not enough.

{\small
\begin{longtable}{@{}>{\footnotesize\raggedright\arraybackslash}p{4.2cm}@{\hspace{6pt}}>{\raggedright\arraybackslash}p{10.4cm}@{}}
\toprule
\textbf{Command} & \textbf{What it shows} \\
\midrule
\endhead
\bottomrule
\endlastfoot
\cmd{device} & The device entries that could be added, with their friendly names and providers. \\
\cmd{device activate <name>} & Makes another device the active one, so its settings are the ones in force. \\
\cmd{devinfo} & The current logical and native position, printed in the form \cmd{x,y;nx,ny;}. Compare the two to see how the driver's transform is placing the scene on the bed. \\
\end{longtable}
}

Settings can also be read and written from the console. The \cmd{set} command
works on variables in a context: \cmd{set -p <context> <key> <value>} writes one
value, and \cmd{set -p <context>} with no value lists what that context
currently holds --- the reliable way to confirm the exact key names and their
present values on your installation.

For diagnosis, the driver's own controller window shows the packet flow and the
counters, the \panel{Buffer} window shows the command buffer being fed to the
device, and the \ui{UsbConnect} window shows the USB pipe and accepts typed
input. What those windows report, and what to do about it, is covered in
Chapter~\ref{ch:connect}.
